For a prescribed force $f$ in the linearized Euler equation
\eqref{cc:intro:linearized-momentum}, we choose the initial value $b(0)$
of the phase-independent correction to preserve each particle's final
position to first order. The first-order displacement $\xi$ induced by
$b$ vanishes initially because the starting positions are fixed.
We impose $\xi(T)=0$ to select $b(0)$.

Fix $0<T\le1$ and $I=[0,T]$. Let the parent $u$ be a smooth
divergence-free field on $I\times\mathbb T^3$, not assumed to solve any
equation, and let $X(t,a)$, $X(0,a)=a$, be its volume-preserving flow. Set
\begin{equation}\label{cc:uf:mean:data}
 M=(\nabla u)\circ X,\quad
 H^\#=-(M_t+M^2),\quad
 K=(D_aX)^TD_aX,\quad \Omega=(D_aX)^T(M^T-M)D_aX.
\end{equation}
Time derivatives keep $a$ fixed. Choose
$D_-\ge\max\{1,\sup_{I,a}|(D_aX)^{-1}|\}$.
The flow equation gives
\begin{equation}\label{cc:uf:mean:flow-identities}
 \partial_t D_aX=MD_aX,\qquad
 K_t-\Omega=2(D_aX)^T\partial_t D_aX.
\end{equation}
Let $p$ be a smooth function and write $R=u_t+u\cdot\nabla u+\nabla p$.
The gradient of this identity is
$(\partial_t+u\cdot\nabla)\nabla u+(\nabla u)^2+D^2p=\nabla R$, with
$(\nabla u)_{ij}=\partial_ju_i$. Along trajectories the material
derivative becomes $\partial_t$, so
\begin{equation}\label{cc:uf:mean:eq:Hsharp}
 H^\#=(D^2p-\nabla R)\circ X.
\end{equation}
For a solution of the Euler equations, $H^\#$ is the pressure Hessian
along the flow: it is what prevents the velocity gradient from evolving
by the Riccati equation $M_t=-M^2$, and it is symmetric; for a linear
flow $Dx$ this is the condition that $D'+D^2$ be symmetric in
Subsection~\ref{cc:sec:affine}. In general $H^\#$ need not be symmetric.
The matrix $M^T-M$ is the vorticity of $u$ along the flow:
$(M^T-M)y=-(\curl u\circ X)\times y$. Thus $\Omega$ is the vorticity
pulled back to labels. When $H^\#\le\Id$ as quadratic forms, the
displacement energy is coercive and the inverse in Appendix~\ref{cc:app:mean}
gives a unique correction with a H\"older bound.

\subsection{Displacement and the two end conditions}

Here $H^1_0(I;\R^3)$ is the space of $z\in H^1(I;\R^3)$ with
$z(0)=z(T)=0$, the closure of $C^\infty_c((0,T);\R^3)$ in $H^1$. Fix a
label $a$, let $z\in H^1_0(I;\R^3)$, and put $\xi=D_aX(\cdot,a)z$. Since
$\partial_t D_aX=MD_aX$,
\begin{equation}\label{cc:uf:mean:xi-t}
 \xi_t=(\partial_t D_aX)z+D_aX z_t=M\xi+D_aX z_t,\qquad
 z_t=(D_aX)^{-1}(\xi_t-M\xi).
\end{equation}

By the first variation formula \eqref{cc:intro:first-variation}, the
displacement induced by $b$ satisfies $\xi_t=M\xi+b\circ X$.
Thus \eqref{cc:uf:mean:xi-t} identifies $D_aX z_t=b\circ X$.
Since $D_aX$ is invertible, the prescribed final displacement
$\xi(T)=0$ is equivalent to $z(T)=0$.

Cauchy--Schwarz in $z(t)=-\int_t^T z_t(s)\dd s$ gives
\begin{equation}\label{cc:uf:mean:poincare}
 |z(t)|^2\le(T-t)\int_I|z_t|^2\dd s,\qquad
 \int_I|z|^2\dd t\le\frac{T^2}{2}\int_I|z_t|^2\dd t.
\end{equation}
Since $D_aX$ is smooth, $\xi\in H^1_0(I;\R^3)$ as well, so
\eqref{cc:uf:mean:poincare} also holds for $\xi$; in particular
$\|\xi\|_{L^2(I)}\le\|\xi_t\|_{L^2(I)}$, since $T\le1$. By
\eqref{cc:uf:mean:xi-t} and $|(D_aX)^{-1}|\le D_-$,
\[
 \|z_t\|_{L^2(I)}
 \le D_-\bigl(\|\xi_t\|_{L^2(I)}+\|M\|_\infty\|\xi\|_{L^2(I)}\bigr)
 \le D_-(1+\|M\|_\infty)\|\xi_t\|_{L^2(I)}.
\]
Squaring gives
\begin{equation}\label{cc:uf:mean:equivalence}
 \int_I|z_t|^2\dd t\le D_-^2(1+\|M\|_\infty)^2\int_I|\xi_t|^2\dd t.
\end{equation}

\subsection{Coercivity}

The inverse in Appendix~\ref{cc:app:mean} requires the label form
$\int_I(Kz_t\cdot z_t+\Omega z_t\cdot z)\dd t$ to control
$\int_I|z_t|^2\dd t$. In the displacement variable $\xi$, the integrand
contains $H^\#$ and a total derivative that integrates to zero because
$\xi$ vanishes at both ends. We can then use the upper bound on $H^\#$ and
\eqref{cc:uf:mean:equivalence} to obtain the required estimate.

Substitute $D_aX z_t=\xi_t-M\xi$ and $D_aX z=\xi$ in the two terms. First,
$Kz_t\cdot z_t=|D_aX z_t|^2=|\xi_t|^2-2\xi_t\cdot M\xi+|M\xi|^2$. Second,
$\Omega z_t\cdot z=(M^T-M)(\xi_t-M\xi)\cdot\xi
=\xi_t\cdot M\xi-\xi_t\cdot M^T\xi-|M\xi|^2+M^2\xi\cdot\xi$. Adding,
and using
$\partial_t(\xi\cdot M\xi)=\xi_t\cdot M\xi+\xi_t\cdot M^T\xi+M_t\xi\cdot\xi$
and $M_t+M^2=-H^\#$, we obtain
\[
 Kz_t\cdot z_t+\Omega z_t\cdot z
 =|\xi_t|^2-H^\#\xi\cdot\xi-\partial_t(\xi\cdot M\xi).
\]
Since $\xi$ vanishes at both ends, integration and density give
\begin{equation}\label{cc:uf:mean:jacobi-id}
 \int_I(Kz_t\cdot z_t+\Omega z_t\cdot z)\dd t
 =\int_I(|\xi_t|^2-H^\#\xi\cdot\xi)\dd t.
\end{equation}

\begin{lemma}\label{cc:uf:mean:hardy}
Assume $H^\#\le\Id$ on $I\times\mathbb T^3$ as quadratic forms. Then
for every $a\in\mathbb T^3$, $z\in H^1_0(I;\R^3)$ and
$\xi=D_aX(\cdot,a)z$, one has
$\int_I H^\#\xi\cdot\xi\dd t\le\frac12\int_I|\xi_t|^2\dd t$. Hence
\begin{equation}\label{cc:uf:mean:coercive}
 \int_I(Kz_t\cdot z_t+\Omega z_t\cdot z)\dd t
 \ge\frac12\int_I|\xi_t|^2\dd t.
\end{equation}
\end{lemma}
\begin{proof}
The assumption $H^\#\le\Id$, the second inequality of
\eqref{cc:uf:mean:poincare} for $\xi$ and $T\le1$ give
\[
 \int_I H^\#\xi\cdot\xi\dd t
 \le\int_I|\xi|^2\dd t
 \le\frac{T^2}2\int_I|\xi_t|^2\dd t
 \le\frac12\int_I|\xi_t|^2\dd t.
\]
Equation~\eqref{cc:uf:mean:jacobi-id} gives
\eqref{cc:uf:mean:coercive}.
\end{proof}

\subsection{The inverse and its pressure}\label{cc:uf:sec:mean}

Given a source $L(t,a)$, we solve the system
\begin{equation}\label{cc:uf:mean:integrated}
 K(t)v(t)-\int_0^t(\Omega v)(s)\dd s+L(t)=c+\nabla_a\phi,
 \qquad v=z_t,\quad z(0)=z(T)=0,\quad\diver_a z=0,
\end{equation}
for the displacement $z$, a time-independent divergence-free field $c$
and a label pressure gradient $\nabla_a\phi$. For the force sources of
Theorem~\ref{cc:uf:mean:inverse}, \eqref{cc:uf:mean:integrated} is the
linearized momentum equation \eqref{cc:uf:mean:strong} integrated once
in time. For these sources $L(0)=0$. Since $K(0)=\Id$, the equation
at time zero reads $v(0)=c+\nabla_a\phi(0)$. Projecting
onto divergence-free fields gives $c=v(0)$. Since $b\circ X=D_aXv$
and $X(0,a)=a$, the constant $c$ is the initial correction $b(0)$.

\begin{theorem}\label{cc:uf:mean:inverse}
Assume $H^\#\le\Id$ on $I\times\mathbb T^3$ as quadratic forms. For
every $L\in C(I;\Holder{0})$ the system
\eqref{cc:uf:mean:integrated} has a unique solution with
$z\in C^1(I;\Holder{0})$, $c\in\Holder{0}$ divergence free and
$\nabla_a\phi\in C(I;\Holder{0})$. There are absolute $C,p$ such that, with
\begin{equation}\label{cc:uf:mean:base-constant}
 A=C\bigl(D_-+\|M\|_\infty+\norm K_{C_t\Holder{1}}
   +\norm\Omega_{C_t\Holder{1}}\bigr)^p,
\end{equation}
one has $\norm v_{C_t\Holder{0}}\le A\norm L_{C_t\Holder{0}}$.

Let $f$ be a smooth field on $I\times\mathbb T^3$, and let
$L=-\int_0^tg(s)\dd s$ with $g=(D_aX)^T(f\circ X)$. Then $z$ is smooth.
Put $B=D_aX v$ and $\mathcal D=2(D_aX)^T\partial_t D_aX$.
There is a smooth zero-mean label pressure $\mathsf q$ such that
\begin{equation}\label{cc:uf:mean:strong}
 Kv_t+\nabla_a\mathsf q=g-\mathcal Dv.
\end{equation}
Thus $b=B\circ X^{-1}$ and $q=\mathsf q\circ X^{-1}$ are smooth and solve
\begin{equation}\label{cc:uf:mean:physical}
 \diver b=0,\qquad
 b_t+u\cdot\nabla b+b\cdot\nabla u+\nabla q=f.
\end{equation}
The endpoint conditions select the initial velocity $b(0)$.

The number $A$ also bounds, uniformly in $t$, the operator norm on
$\Holder{0}$ of the fixed-time inverse $(\PP K(t))^{-1}\PP$, where
$(\PP K(t))^{-1}$ inverts the map $v\mapsto\PP(K(t)v)$ on the
divergence-free fields of $\Holder{0}$.
\end{theorem}
\begin{proof}
Rescale time by $\tau=t/T$ and put $Z(\tau)=z(t)/T$, so
$Z_\tau=v$. Theorem~\ref{cc:holder:linear} applies with coefficients
$K,T\Omega$, ellipticity $D_-^{-2}$ and source $L$.
Its coercivity constant is $1/\bigl(2D_-^2(1+\|M\|_\infty)^2\bigr)$,
by Lemma~\ref{cc:uf:mean:hardy} and \eqref{cc:uf:mean:equivalence}.
It gives the unique solution. Since $D_-\ge1$ and $T\le1$, the quantity
$Q$ of Appendix~\ref{cc:app:mean} is at most an absolute constant times
a fixed power of the sum in parentheses in
\eqref{cc:uf:mean:base-constant}, so the estimate
\eqref{cc:holder:eq:base-linear} gives the asserted bound.
For a smooth $f$ the source $L$ is smooth, and so is the solution, by
the last assertion of Theorem~\ref{cc:holder:linear}.
Differentiating \eqref{cc:uf:mean:integrated}, using
$K_t-\Omega=\mathcal D$ from \eqref{cc:uf:mean:flow-identities}, and
setting $\mathsf q=-\phi_t$ gives
\eqref{cc:uf:mean:strong}. Then
$B_t+MB=D_aX v_t+2(\partial_t D_aX)v$ and the Piola identity
$\diver_x(B\circ X^{-1})=(\diver_a v)\circ X^{-1}$ give
\eqref{cc:uf:mean:physical}.
Finally, by Lemma~\ref{cc:holder:freezing} with $H=\R^3$, $\mathcal R=\PP$ and
$B=K(t)$, whose coercivity is $D_-^{-2}$, the fixed-time inverse is
bounded by a fixed power of $2+D_-^2+\|K(t)\|_{C^1}$. Since
$D_-\ge1$, this number is also at most an absolute constant times a
fixed power of the sum in parentheses in
\eqref{cc:uf:mean:base-constant}, so $A$ bounds the fixed-time inverse
once $C$ and $p$ are enlarged.
\end{proof}
